\documentclass[12pt,a4paper]{article}
\usepackage[utf8]{inputenc}
\usepackage[T5]{fontenc}
\usepackage{amsmath, amssymb}
\usepackage{graphicx}
\usepackage{ifthen}

\newboolean{showsolutions}
\setboolean{showsolutions}{true}

% --- ĐỊNH NGHĨA CÁC LỆNH ẢO CHO CÔNG CỤ LATEX2JSON CỦA WEBSITE ---
\newcommand{\doctitle}[1]{\begin{center}\Large\textbf{#1}\end{center}}
\newcommand{\docdesc}[1]{\textbf{Mô tả:} #1}
\newcommand{\docgrade}[1]{\textbf{Lớp:} #1}
\newcommand{\docstatus}[1]{\textbf{Trạng thái:} #1}
\newcommand{\doctype}[1]{\textbf{Loại:} #1}
\newcommand{\doctopics}[1]{\textbf{Chủ đề:} #1}

\newenvironment{textblock}{}{}

\newenvironment{lesson}[2]{
	\noindent\textbf{\Large #1}\\
	\textit{#2}
}{}

\newcommand{\image}[3]{
	\begin{center}
		\includegraphics[width=#1\textwidth]{#3} \\
		\textit{#2}
	\end{center}
}

\newenvironment{quiz}[1]{
	\noindent\textbf{\Large Kiểm tra: #1}
}{}

\newenvironment{mcq}[4]{
	\noindent\textbf{Câu hỏi:} #1 \\
	\ifthenelse{\boolean{showsolutions}}{
		\textit{(Đáp án đúng: #2, Điểm: #3) - Giải thích: #4}
	}{}
	\begin{itemize}
	}{
	\end{itemize}
}
\newcommand{\option}[1]{\item #1}

\newenvironment{truefalse}[3]{
	\noindent\textbf{Đúng/Sai:} #1 \\
	\ifthenelse{\boolean{showsolutions}}{
		\textit{(Điểm: #2) - Giải thích: #3}
	}{}
	\begin{itemize}
	}{
	\end{itemize}
}
\newcommand{\statement}[2]{\item [\textbf{#1}] #2}

\newcommand{\shortanswer}[4]{
    \noindent\textbf{Trả lời ngắn (Điểm: #3):}

    \noindent #1

	\ifthenelse{\boolean{showsolutions}}{
		\noindent\textit{Đáp án: #2 - Giải thích:}
		\noindent #4
	}{}
}

\newcommand{\essay}[3]{
    \noindent\textbf{Tự luận (Điểm: #2):}
    
    \noindent #1
    
	\ifthenelse{\boolean{showsolutions}}{
		\noindent\textbf{Gợi ý / Đáp án:}
		\noindent #3
	}{}
}

\begin{document}

\doctitle{ĐỀ SỐ 2: BÀI TẬP RÈN LUYỆN PHƯƠNG TRÌNH MẶT CẦU}
\docdesc{Đề kiểm tra rèn luyện số 2 chuyên đề Phương trình mặt cầu - Toán 12 SGK Mới}
\docgrade{12}
\docstatus{draft}
\doctype{test}
\doctopics{Hình học không gian, Oxyz, Phương trình mặt cầu}

\begin{quiz}{Phần I: Trắc nghiệm nhiều phương án lựa chọn}

\begin{mcq}{Cho điểm $M$ nằm ngoài mặt cầu $S(O;R)$. Khẳng định nào dưới đây đúng?}{2}{1}{Điểm $M$ nằm ngoài mặt cầu $S(O;R) \Leftrightarrow OM > R$.

Chọn C.}
	\option{$OM < R$.}
	\option{$OM = R$.}
	\option{$OM > R$.}
	\option{$OM \le R$.}
\end{mcq}

\begin{mcq}{Trong không gian $Oxyz$, cho mặt cầu $(S): x^2 + (y-2)^2 + (z+1)^2 = 6$. Đường kính của $(S)$ bằng}{2}{1}{Bán kính mặt cầu là $R = \sqrt{6}$.
Độ dài đường kính bằng $2R = 2\sqrt{6}$.

Chọn C.}
	\option{$\sqrt{6}$.}
	\option{$12$.}
	\option{$2\sqrt{6}$.}
	\option{$3$.}
\end{mcq}

\begin{mcq}{Mặt cầu $(S): 3x^2 + 3y^2 + 3z^2 - 6x + 12y + 2 = 0$ có bán kính bằng}{3}{1}{Chia hai vế cho 3: $x^2 + y^2 + z^2 - 2x + 4y + \frac{2}{3} = 0$.
Ta có $a = 1, b = -2, c = 0, d = \frac{2}{3}$.
Bán kính $R = \sqrt{1^2 + (-2)^2 + 0^2 - \frac{2}{3}} = \sqrt{\frac{13}{3}}$.

Chọn D.}
	\option{$\frac{\sqrt{7}}{3}$.}
	\option{$\frac{2\sqrt{7}}{3}$.}
	\option{$\frac{\sqrt{21}}{3}$.}
	\option{$\sqrt{\frac{13}{3}}$.}
\end{mcq}

\begin{mcq}{Trong không gian $Oxyz$, cho mặt cầu $(S): (x-2)^2 + (y+1)^2 + (z-3)^2 = 4$. Tâm của $(S)$ có tọa độ là}{3}{1}{Tâm của mặt cầu $(S)$ là $I(2;-1;3)$.

Chọn D.}
	\option{$(-2;1;-3)$.}
	\option{$(-4;2;-6)$.}
	\option{$(4;-2;6)$.}
	\option{$(2;-1;3)$.}
\end{mcq}

\begin{mcq}{Trong không gian $Oxyz$, mặt cầu $(S): (x+1)^2 + (y-2)^2 + z^2 = 9$ có bán kính bằng}{0}{1}{Bán kính của mặt cầu là $R = \sqrt{9} = 3$.

Chọn A.}
	\option{$3$.}
	\option{$81$.}
	\option{$9$.}
	\option{$6$.}
\end{mcq}

\begin{mcq}{Trong không gian $Oxyz$, cho mặt cầu $(S)$ có tâm $I(1;-4;0)$ và bán kính bằng $3$. Phương trình của $(S)$ là}{1}{1}{Phương trình mặt cầu $(S): (x-1)^2 + (y+4)^2 + z^2 = 3^2 = 9$.

Chọn B.}
	\option{$(x+1)^2 + (y-4)^2 + z^2 = 9$.}
	\option{$(x-1)^2 + (y+4)^2 + z^2 = 9$.}
	\option{$(x-1)^2 + (y+4)^2 + z^2 = 3$.}
	\option{$(x+1)^2 + (y-4)^2 + z^2 = 3$.}
\end{mcq}

\begin{mcq}{Trong không gian $Oxyz$ cho điểm $I(2;3;4)$ và $A(1;2;3)$. Phương trình mặt cầu tâm $I$ và đi qua $A$ có phương trình là:}{3}{1}{Bán kính $R = IA = \sqrt{(1-2)^2 + (2-3)^2 + (3-4)^2} = \sqrt{3} \Rightarrow R^2 = 3$.
Phương trình mặt cầu $(S): (x-2)^2 + (y-3)^2 + (z-4)^2 = 3$.

Chọn D.}
	\option{$(x+2)^2 + (y+3)^2 + (z+4)^2 = 3$.}
	\option{$(x+2)^2 + (y+3)^2 + (z+4)^2 = 9$.}
	\option{$(x-2)^2 + (y-3)^2 + (z-4)^2 = 45$.}
	\option{$(x-2)^2 + (y-3)^2 + (z-4)^2 = 3$.}
\end{mcq}

\begin{mcq}{Trong không gian với hệ trục tọa độ $Oxyz$, cho hai điểm $A(1;2;3), B(5;4;-1)$. Phương trình mặt cầu đường kính $AB$ là}{0}{1}{Tâm $I$ là trung điểm $AB \Rightarrow I(3;3;1)$.
Bán kính $R = \frac{AB}{2} = \frac{\sqrt{(5-1)^2 + (4-2)^2 + (-1-3)^2}}{2} = \frac{\sqrt{36}}{2} = 3 \Rightarrow R^2 = 9$.
Phương trình mặt cầu $(S): (x-3)^2 + (y-3)^2 + (z-1)^2 = 9$.

Chọn A.}
	\option{$(x-3)^2 + (y-3)^2 + (z-1)^2 = 9$.}
	\option{$(x-3)^2 + (y-3)^2 + (z-1)^2 = 6$.}
	\option{$(x+3)^2 + (y+3)^2 + (z+1)^2 = 9$.}
	\option{$(x-3)^2 + (y-3)^2 + (z-1)^2 = 36$.}
\end{mcq}

\begin{mcq}{Trong không gian $Oxyz$, cho hai điểm $M(3;-2;5), N(-1;6;-3)$. Mặt cầu đường kính $MN$ có phương trình là:}{3}{1}{Tâm $I$ là trung điểm $MN \Rightarrow I(1;2;1)$.
Bán kính $R = \frac{MN}{2} = \frac{\sqrt{(-1-3)^2 + (6+2)^2 + (-3-5)^2}}{2} = \frac{\sqrt{144}}{2} = 6 \Rightarrow R^2 = 36$.
Phương trình mặt cầu: $(x-1)^2 + (y-2)^2 + (z-1)^2 = 36$.

Chọn D.}
	\option{$(x+1)^2 + (y+2)^2 + (z+1)^2 = 6$.}
	\option{$(x-1)^2 + (y-2)^2 + (z-1)^2 = 6$.}
	\option{$(x+1)^2 + (y+2)^2 + (z+1)^2 = 36$.}
	\option{$(x-1)^2 + (y-2)^2 + (z-1)^2 = 36$.}
\end{mcq}

\begin{mcq}{Cho hai điểm $A, B$ cố định trong không gian có độ dài $AB$ là $4$. Biết rằng tập hợp các điểm $M$ trong không gian sao cho $MA = 3MB$ là một mặt cầu. Bán kính mặt cầu đó bằng}{3}{1}{Gọi $E, F$ lần lượt là các điểm chia trong và chia ngoài của đoạn thẳng $AB$ theo tỉ số $k = 3$.
Ta có $EA = 3EB \Rightarrow EB = \frac{AB}{4} = 1$.
$FA = 3FB \Rightarrow FA - FB = 2FB = AB = 4 \Rightarrow FB = 2$.
Đường kính của mặt cầu là $EF = EB + BF = 1 + 2 = 3$.
Bán kính mặt cầu là $R = \frac{EF}{2} = \frac{3}{2}$.

Chọn D.}
	\option{$3$.}
	\option{$\frac{9}{2}$.}
	\option{$1$.}
	\option{$\frac{3}{2}$.}
\end{mcq}

\begin{mcq}{Trong không gian với hệ tọa độ $Oxyz$, mặt cầu $(S)$ qua bốn điểm $A(3;3;0), B(3;0;3), C(0;3;3), D(3;3;3)$. Phương trình mặt cầu $(S)$ là}{3}{1}{Tâm mặt cầu $I\left(\frac{3}{2}; \frac{3}{2}; \frac{3}{2}\right)$.
Bán kính $R = ID = \sqrt{\left(3-\frac{3}{2}\right)^2 + \left(3-\frac{3}{2}\right)^2 + \left(3-\frac{3}{2}\right)^2} = \sqrt{\frac{27}{4}}$.
Phương trình mặt cầu $(S): \left(x-\frac{3}{2}\right)^2 + \left(y-\frac{3}{2}\right)^2 + \left(z-\frac{3}{2}\right)^2 = \frac{27}{4}$.

Chọn D.}
	\option{$\left(x-\frac{3}{2}\right)^2 + \left(y-\frac{3}{2}\right)^2 + \left(z-\frac{3}{2}\right)^2 = \frac{3\sqrt{3}}{2}$.}
	\option{$\left(x-\frac{3}{2}\right)^2 + \left(y+\frac{3}{2}\right)^2 + \left(z-\frac{3}{2}\right)^2 = \frac{27}{4}$.}
	\option{$\left(x-\frac{3}{2}\right)^2 + \left(y-\frac{3}{2}\right)^2 + \left(z+\frac{3}{2}\right)^2 = \frac{27}{4}$.}
	\option{$\left(x-\frac{3}{2}\right)^2 + \left(y-\frac{3}{2}\right)^2 + \left(z-\frac{3}{2}\right)^2 = \frac{27}{4}$.}
\end{mcq}

\begin{mcq}{Trong không gian $Oxyz$, cho mặt phẳng $(P): 2x + 3y + z - 11 = 0$. Mặt cầu $(S)$ có tâm $I(1;-2;1)$ và tiếp xúc với mặt phẳng $(P)$ tại điểm $H$, khi đó $H$ có tọa độ là}{3}{1}{Đường thẳng $IH$ qua $I(1;-2;1)$ và vuông góc với $(P)$ nên nhận $\vec{n}=(2;3;1)$ làm VTCP: $\begin{cases} x = 1+2t \\ y = -2+3t \\ z = 1+t \end{cases}$.
Thay vào phương trình mặt phẳng $(P)$:
$2(1+2t) + 3(-2+3t) + (1+t) - 11 = 0 \Leftrightarrow 14t - 14 = 0 \Leftrightarrow t = 1$.
Suy ra $H(3;1;2)$.

Chọn D.}
	\option{$H(-3;-1;-2)$.}
	\option{$H(-1;-5;0)$.}
	\option{$H(1;5;0)$.}
	\option{$H(3;1;2)$.}
\end{mcq}

\begin{mcq}{Trong không gian với hệ trục tọa độ $Oxyz$ cho mặt phẳng $(\alpha)$ có phương trình $2x + y - z - 1 = 0$ và mặt cầu $(S)$ có phương trình $(x-1)^2 + (y-1)^2 + (z+2)^2 = 4$. Xác định bán kính $r$ của đường tròn là giao tuyến của mặt phẳng $(\alpha)$ và mặt cầu $(S)$.}{1}{1}{Tâm $I(1;1;-2)$, bán kính $R = 2$.
Khoảng cách $d(I, (\alpha)) = \frac{|2.1 + 1 - (-2) - 1|}{\sqrt{2^2 + 1^2 + (-1)^2}} = \frac{4}{\sqrt{6}}$.
Bán kính đường tròn giao tuyến là $r = \sqrt{R^2 - d^2} = \sqrt{4 - \frac{16}{6}} = \sqrt{\frac{8}{6}} = \frac{2\sqrt{3}}{3}$.

Chọn B.}
	\option{$r = \frac{2\sqrt{42}}{3}$.}
	\option{$r = \frac{2\sqrt{3}}{3}$.}
	\option{$r = \frac{2\sqrt{15}}{3}$.}
	\option{$r = \frac{2\sqrt{7}}{3}$.}
\end{mcq}

\begin{mcq}{Trong không gian với hệ tọa độ $Oxyz$, cho điểm $I(-3;0;1)$. Mặt cầu $(S)$ có tâm $I$ và cắt mặt phẳng $(P): x - 2y - 2z - 1 = 0$ theo một thiết diện là một hình tròn. Diện tích của hình tròn này bằng $\pi$. Phương trình mặt cầu $(S)$ là}{2}{1}{Bán kính đường tròn giao tuyến $r = \sqrt{\frac{S}{\pi}} = 1$.
Khoảng cách $d(I, (P)) = \frac{|-3 - 2.0 - 2.1 - 1|}{\sqrt{1^2 + (-2)^2 + (-2)^2}} = \frac{6}{3} = 2$.
Bán kính mặt cầu $R = \sqrt{r^2 + d^2} = \sqrt{1^2 + 2^2} = \sqrt{5} \Rightarrow R^2 = 5$.
Phương trình mặt cầu $(S): (x+3)^2 + y^2 + (z-1)^2 = 5$.

Chọn C.}
	\option{$(x+3)^2 + y^2 + (z-1)^2 = 4$.}
	\option{$(x+3)^2 + y^2 + (z-1)^2 = 25$.}
	\option{$(x+3)^2 + y^2 + (z-1)^2 = 5$.}
	\option{$(x+3)^2 + y^2 + (z-1)^2 = 2$.}
\end{mcq}

\begin{mcq}{Trong không gian với hệ tọa độ $Oxyz$, cho mặt phẳng $(P): x - 2y + 2z - 2 = 0$ và điểm $I(-1; 2; -1)$. Viết phương trình mặt cầu $(S)$ có tâm $I$ và cắt mặt phẳng $(P)$ theo giao tuyến là đường tròn có bán kính bằng $5$.}{3}{1}{Khoảng cách $d(I, (P)) = \frac{|-1 - 2.2 + 2.(-1) - 2|}{\sqrt{1^2 + (-2)^2 + 2^2}} = \frac{9}{3} = 3$.
Bán kính mặt cầu $R = \sqrt{r^2 + d^2} = \sqrt{5^2 + 3^2} = \sqrt{34} \Rightarrow R^2 = 34$.
Phương trình mặt cầu $(S): (x+1)^2 + (y-2)^2 + (z+1)^2 = 34$.

Chọn D.}
	\option{$(S): (x+1)^2 + (y-2)^2 + (z+1)^2 = 25$.}
	\option{$(S): (x+1)^2 + (y-2)^2 + (z+1)^2 = 16$.}
	\option{$(S): (x-1)^2 + (y+2)^2 + (z-1)^2 = 34$.}
	\option{$(S): (x+1)^2 + (y-2)^2 + (z+1)^2 = 34$.}
\end{mcq}

\begin{mcq}{Trong không gian $Oxyz$, mặt phẳng $(P): 2x + 6y + z - 3 = 0$ cắt trục $Oz$ và đường thẳng $d: \frac{x-5}{1} = \frac{y}{2} = \frac{z-6}{-1}$ lần lượt tại $A$ và $B$. Phương trình mặt cầu đường kính $AB$ là}{1}{1}{$A = (P) \cap Oz \Rightarrow A(0;0;3)$.
$B = (P) \cap d \Rightarrow B(4;-2;7)$.
Tâm $I$ là trung điểm $AB \Rightarrow I(2;-1;5)$.
Bán kính $R = \frac{AB}{2} = \frac{\sqrt{4^2 + (-2)^2 + 4^2}}{2} = 3 \Rightarrow R^2 = 9$.
Phương trình $(S): (x-2)^2 + (y+1)^2 + (z-5)^2 = 9$.

Chọn B.}
	\option{$(x+2)^2 + (y-1)^2 + (z+5)^2 = 36$.}
	\option{$(x-2)^2 + (y+1)^2 + (z-5)^2 = 9$.}
	\option{$(x+2)^2 + (y-1)^2 + (z+5)^2 = 9$.}
	\option{$(x-2)^2 + (y+1)^2 + (z-5)^2 = 36$.}
\end{mcq}

\begin{mcq}{Trong không gian $Oxyz$, cho đường thẳng $d: \frac{x}{2} = \frac{y-3}{1} = \frac{z-2}{1}$ và hai mặt phẳng $(P): x - 2y + 2z = 0; (Q): x - 2y + 3z - 5 = 0$. Mặt cầu $(S)$ có tâm $I$ là giao điểm của đường thẳng $(d)$ và mặt phẳng $(P)$. Mặt phẳng $(Q)$ tiếp xúc với mặt cầu $(S)$. Viết phương trình mặt cầu $(S)$.}{2}{1}{$I = d \cap (P) \Rightarrow I(2;4;3)$.
Bán kính $R = d(I, (Q)) = \frac{|2 - 2.4 + 3.3 - 5|}{\sqrt{1^2 + (-2)^2 + 3^2}} = \frac{2}{\sqrt{14}} \Rightarrow R^2 = \frac{4}{14} = \frac{2}{7}$.
Phương trình mặt cầu $(S): (x-2)^2 + (y-4)^2 + (z-3)^2 = \frac{2}{7}$.

Chọn C.}
	\option{$(S): (x+2)^2 + (y+4)^2 + (z+3)^2 = 1$.}
	\option{$(S): (x-2)^2 + (y-4)^2 + (z-3)^2 = 6$.}
	\option{$(S): (x-2)^2 + (y-4)^2 + (z-3)^2 = \frac{2}{7}$.}
	\option{$(S): (x-2)^2 + (y+4)^2 + (z+4)^2 = 8$.}
\end{mcq}

\begin{mcq}{Trong không gian $Oxyz$, cho các điểm $I(1;0;0)$ và đường thẳng $d: \frac{x-1}{1} = \frac{y-1}{2} = \frac{z+2}{1}$. Phương trình mặt cầu $(S)$ có tâm $I$ và cắt đường thẳng $d$ tại hai điểm $A, B$ sao cho tam giác $IAB$ đều là}{1}{1}{Tam giác $IAB$ đều cạnh $AB = R$ nên khoảng cách $h = d(I,d) = \frac{R\sqrt{3}}{2} \Rightarrow R^2 = \frac{4h^2}{3}$.
Lấy $M(1;1;-2) \in d \Rightarrow \vec{IM} = (0;1;-2)$. VTCP $\vec{u} = (1;2;1)$.
$[\vec{IM},\vec{u}] = (5;-2;-1) \Rightarrow h = \frac{\sqrt{25+4+1}}{\sqrt{1+4+1}} = \sqrt{5}$.
Suy ra $R^2 = \frac{4.5}{3} = \frac{20}{3}$.
Phương trình mặt cầu $(S): (x-1)^2 + y^2 + z^2 = \frac{20}{3}$.

Chọn B.}
	\option{$(x+1)^2 + y^2 + z^2 = \frac{20}{3}$.}
	\option{$(x-1)^2 + y^2 + z^2 = \frac{20}{3}$.}
	\option{$(x-1)^2 + y^2 + z^2 = \frac{16}{4}$.}
	\option{$(x-1)^2 + y^2 + z^2 = \frac{5}{3}$.}
\end{mcq}

\end{quiz}

\begin{quiz}{Phần II: Trắc nghiệm đúng sai}

\begin{truefalse}{Trong không gian với hệ toạ độ $Oxyz$, cho điểm $M(2;0;2)$ và mặt cầu $(S): x^2 + (y+2)^2 + (z-2)^2 = 8$. Xét tính đúng sai của các khẳng định sau:}{1}{- a) Đúng. Thay $M(2;0;2)$ vào $(S)$: $2^2 + 2^2 + 0^2 = 8$ nên $M \in (S)$.
- b) Đúng. $R = \sqrt{8} = 2\sqrt{2}$.
- c) Đúng. Tâm của mặt cầu $(S)$ là $I(0;-2;2)$.
- d) Sai. Khoảng cách $d(I, (Oxy)) = |z_I| = 2 < R = 2\sqrt{2}$ nên $(Oxy)$ cắt mặt cầu $(S)$ theo đường tròn.}
	\statement{true}{Điểm $M(2;0;2)$ thuộc mặt cầu $(S)$.}
	\statement{true}{Bán kính mặt cầu $(S)$ là $R = 2\sqrt{2}$.}
	\statement{true}{Tọa độ tâm mặt cầu $(S)$ là $I(0;-2;2)$.}
	\statement{false}{Mặt phẳng $(Oxy)$ và mặt cầu $(S)$ không có điểm chung.}
\end{truefalse}

\begin{truefalse}{Trong không gian $Oxyz$, mặt cầu $(S): (x-1)^2 + (y+2)^2 + (z-4)^2 = 20$. Xét tính đúng sai của các khẳng định sau:}{1}{- a) Sai. Bán kính của mặt cầu là $R = \sqrt{20} = 2\sqrt{5}$.
- b) Sai. Tâm của mặt cầu $(S)$ là $I(1;-2;4)$.
- c) Đúng. $d(I, (P)) = \frac{|2.1 + 4 + 4|}{\sqrt{2^2 + 0^2 + 1^2}} = \frac{10}{\sqrt{5}} = 2\sqrt{5} = R \Rightarrow (S)$ tiếp xúc với $(P)$.
- d) Đúng. $d(I, (Ozx)) = |y_I| = |-2| = 2 \Rightarrow r = \sqrt{R^2 - d^2} = \sqrt{20 - 4} = 4$.}
	\statement{false}{Bán kính mặt cầu $(S)$ là 20.}
	\statement{false}{Tọa độ tâm mặt cầu $(S)$ là $I(-1;2;-4)$.}
	\statement{true}{Mặt cầu $(S)$ tiếp xúc với mặt phẳng $(P): 2x + z + 4 = 0$.}
	\statement{true}{Mặt phẳng $(Ozx)$ cắt mặt cầu $(S)$ theo đường tròn giao tuyến có bán kính $r = 4$.}
\end{truefalse}

\begin{truefalse}{Trong không gian với hệ tọa độ $Oxyz$, cho các phương trình sau:
$(S_1): x^2 + y^2 + z^2 + x - 2y + 4z - 3 = 0$;
$(S_2): 2x^2 + 2y^2 + 2z^2 - x - y - z = 0$;
$(S_3): 2x^2 + 2y^2 + 2z^2 + 4x + 8y + 6z + 3 = 0$;
$(S_4): x^2 + y^2 + z^2 - 2x + 4y - 4z + 10 = 0$.
Xét tính đúng sai của các khẳng định sau:}{1}{- a) Đúng. $(S_1)$ có $a^2+b^2+c^2-d = \frac{1}{4} + 1 + 4 - (-3) = \frac{33}{4} > 0$.
- b) Đúng. $(S_2)$ đưa về $x^2+y^2+z^2 - \frac{1}{2}x - \frac{1}{2}y - \frac{1}{2}z = 0$ có $a^2+b^2+c^2-d = \frac{3}{16} > 0$.
- c) Sai. $(S_3)$ đưa về $x^2+y^2+z^2+2x+4y+3z+\frac{3}{2}=0$ có $a^2+b^2+c^2-d = 1+4+\frac{9}{4}-\frac{3}{2} = \frac{23}{4} > 0$ nên $(S_3)$ là phương trình mặt cầu.
- d) Đúng. $(S_4)$ có $a^2+b^2+c^2-d = 1+4+4-10 = -1 < 0$ nên không phải là phương trình mặt cầu.}
	\statement{true}{$(S_1)$ là phương trình của một mặt cầu.}
	\statement{true}{$(S_2)$ là phương trình của một mặt cầu.}
	\statement{false}{$(S_3)$ không phải là phương trình của một mặt cầu.}
	\statement{true}{$(S_4)$ không phải là phương trình của một mặt cầu.}
\end{truefalse}

\begin{truefalse}{Trong không gian với hệ trục tọa độ $Oxyz$, cho mặt cầu $(S)$ có tâm nằm trên mặt phẳng $Oxy$ và đi qua ba điểm $A(1;2;-4), B(1;-3;1), C(2;2;3)$. Xét tính đúng sai của các khẳng định sau:}{1}{- a) Sai. Tọa độ tâm là $I(-2;1;0)$.
- b) Đúng. Bán kính $R = IA = \sqrt{(1+2)^2 + (2-1)^2 + (-4)^2} = \sqrt{26}$.
- c) Đúng. Giao với $Ox (y=0, z=0)$: $x^2 + 4x - 21 = 0 \Leftrightarrow x = 3$ hoặc $x = -7$ (hai nghiệm trái dấu).
- d) Sai. Bán kính đường tròn ngoại tiếp $\triangle ABC$ là $r = \frac{5\sqrt{170}}{17} \approx 3{,}83 > 3$.}
	\statement{false}{Tọa độ tâm $I$ của mặt cầu $(S)$ là $I(2;-1;0)$.}
	\statement{true}{Bán kính $R$ của mặt cầu $(S)$ là $R = \sqrt{26}$.}
	\statement{true}{Mặt cầu $(S)$ cắt trục $Ox$ tại hai điểm phân biệt có hoành độ trái dấu.}
	\statement{false}{Mặt phẳng $(ABC)$ cắt mặt cầu theo giao tuyến là đường tròn có bán kính bé hơn 3.}
\end{truefalse}

\begin{truefalse}{Trong không gian với hệ trục tọa độ $Oxyz$, cho điểm $A(1;1;2), B(3;2;-3)$. Mặt cầu $(S)$ có tâm $I$ thuộc $Ox$ và đi qua hai điểm $A, B$. Xét tính đúng sai của các khẳng định sau:}{1}{- a) Đúng. Tâm $I(a;0;0) \in Ox$, từ $IA^2 = IB^2$ giải được $a = 4 \Rightarrow I(4;0;0)$.
- b) Sai. Bán kính $R = IA = \sqrt{(1-4)^2 + 1^2 + 2^2} = \sqrt{14}$.
- c) Đúng. Phương trình mặt cầu $(S): (x-4)^2 + y^2 + z^2 = 14 \Leftrightarrow x^2 + y^2 + z^2 - 8x + 2 = 0$.
- d) Sai. Khoảng cách từ $I$ đến đoạn $AB$ bằng $\frac{\sqrt{26}}{2} \ne \sqrt{26}$.}
	\statement{true}{Tọa độ tâm $(I)$ của mặt cầu $(S)$ là $I(4;0;0)$.}
	\statement{false}{Bán kính $R$ của mặt cầu $(S)$ là $R = 14$.}
	\statement{true}{Mặt cầu $(S)$ có phương trình $x^2 + y^2 + z^2 - 8x + 2 = 0$.}
	\statement{false}{Khoảng cách từ tâm $I$ đến đoạn thẳng $AB$ bằng $\sqrt{26}$.}
\end{truefalse}

\begin{truefalse}{Trong không gian $Oxyz$, mặt cầu $(S)$ đi qua điểm $A(1;-1;4)$ và tiếp xúc với các mặt phẳng tọa độ. Xét tính đúng sai của các khẳng định sau:}{1}{- a) Đúng. Mặt cầu tiếp xúc với 3 mặt phẳng tọa độ nên $|a| = |b| = |c| = R$.
- b) Đúng. Do $A(1;-1;4)$ nên tâm $I(R;-R;R)$. Thay $A$ vào giải được $R = 3 \Rightarrow (x-3)^2 + (y+3)^2 + (z-3)^2 = 9$.
- c) Đúng. Trung điểm $AB$ là $I(3;-3;3)$ và độ dài $AB = 6 = 2R$ nên $AB$ là đường kính.
- d) Sai. Cạnh hình lập phương ngoại tiếp là $2R = 6 \Rightarrow$ đường chéo bằng $6\sqrt{3} \ne 3\sqrt{3}$.}
	\statement{true}{Nếu $I(a;b;c)$ là tâm của mặt cầu $(S)$ thì $|a| = |b| = |c|$.}
	\statement{true}{Mặt cầu $(S)$ có phương trình $(x-3)^2 + (y+3)^2 + (z-3)^2 = 9$.}
	\statement{true}{Mặt cầu $(S)$ đi qua điểm $B(5;-5;2)$ và nhận $AB$ là đường kính.}
	\statement{false}{Mặt cầu $(S)$ nội tiếp một hình lập phương có độ dài đường chéo bằng $3\sqrt{3}$.}
\end{truefalse}

\end{quiz}

\begin{quiz}{Phần III: Trắc nghiệm trả lời ngắn}

\shortanswer{Trong không gian $Oxyz$, cho hai điểm $A(-1;0;1), B(1;-1;2)$. Tập hợp tất cả các điểm $M$ bất kì trong không gian thỏa mãn $\widehat{AMB} = 90^\circ$ là một mặt cầu $(S)$ có bán kính bằng bao nhiêu?}{1{,}2}{1}{Tập hợp điểm $M$ nhìn đoạn $AB$ dưới một góc vuông là mặt cầu đường kính $AB$.
Bán kính $R = \frac{AB}{2} = \frac{\sqrt{(1+1)^2 + (-1-0)^2 + (2-1)^2}}{2} = \frac{\sqrt{6}}{2} \approx 1{,}2$.}

\shortanswer{Trong không gian $Oxyz$, cho hai điểm $A(1;2;1), B(3;1;-2)$. Tập hợp tất cả các điểm $M$ bất kì trong không gian thỏa mãn $MA^2 + MB^2 = 30$ là một mặt cầu $(S)$ có đường kính bằng bao nhiêu?}{6{,}8}{1}{Gọi $I$ là trung điểm $AB \Rightarrow I(2; 1{,}5; -0{,}5)$.
Ta có $MA^2 + MB^2 = 2MI^2 + \frac{AB^2}{2} \Leftrightarrow 30 = 2MI^2 + 7 \Leftrightarrow MI^2 = \frac{23}{2} \Rightarrow R = \sqrt{\frac{23}{2}}$.
Đường kính mặt cầu $d = 2R = \sqrt{46} \approx 6{,}8$.}

\shortanswer{Trong không gian với hệ trục tọa độ $Oxyz$, cho ba điểm $A(1;2;-4), B(1;-3;1), C(2;2;3)$. Tính đường kính của mặt cầu $(S)$ đi qua ba điểm trên và có tâm nằm trên mặt phẳng $(Oxy)$.}{10{,}2}{1}{Tâm mặt cầu là $I(-2;1;0)$, bán kính $R = \sqrt{26}$.
Đường kính của mặt cầu là $d = 2R = 2\sqrt{26} \approx 10{,}2$.}

\shortanswer{Trong không gian $Oxyz$, gọi $(S)$ là mặt cầu đi qua điểm $D(0;1;2)$ và tiếp xúc với các trục $Ox, Oy, Oz$ tại các điểm $A(a;0;0), B(0;b;0), C(0;0;c)$ trong đó $a, b, c \in \mathbb{R} \setminus \{0;1\}$. Tính bán kính của $(S)$.}{7{,}1}{1}{Do $(S)$ tiếp xúc với ba trục tọa độ tại $A(a;0;0), B(0;b;0), C(0;0;c)$ nên tâm $I(a;b;c)$ và $R^2 = a^2+b^2 = b^2+c^2 = c^2+a^2 \Rightarrow a^2 = b^2 = c^2 \Rightarrow R = |a|\sqrt{2}$.
Phương trình $(S): (x-a)^2 + (y-b)^2 + (z-c)^2 = 2a^2$.
$(S)$ đi qua $D(0;1;2) \Rightarrow a^2 + (1-b)^2 + (2-c)^2 = 2a^2 \Leftrightarrow (1-b)^2 + (2-c)^2 = b^2$.
Với $b=c$ ta có $(1-b)^2 + (2-b)^2 = b^2 \Leftrightarrow b^2 - 6b + 5 = 0 \Leftrightarrow b = 5$ (do $b \neq 1$).
Bán kính $R = 5\sqrt{2} \approx 7{,}1$.}

\shortanswer{Trong không gian với hệ trục tọa độ $Oxyz$, cho ba điểm $A(1;0;0), C(0;0;3), B(0;2;0)$. Tập hợp các điểm $M$ thỏa mãn $MA^2 = MB^2 + MC^2$ là mặt cầu có bán kính là bao nhiêu?}{1{,}4}{1}{Gọi $M(x;y;z)$, ta có:
$MA^2 = (x-1)^2 + y^2 + z^2 = x^2 + y^2 + z^2 - 2x + 1$
$MB^2 = x^2 + (y-2)^2 + z^2 = x^2 + y^2 + z^2 - 4y + 4$
$MC^2 = x^2 + y^2 + (z-3)^2 = x^2 + y^2 + z^2 - 6z + 9$
$MA^2 = MB^2 + MC^2 \Leftrightarrow x^2 + y^2 + z^2 + 2x - 4y - 6z + 12 = 0$.
Mặt cầu có tâm $I(-1;2;3)$ và bán kính $R = \sqrt{(-1)^2 + 2^2 + 3^2 - 12} = \sqrt{2} \approx 1{,}4$.}

\shortanswer{Trong không gian $Oxyz$, cho mặt cầu $(S)$ có phương trình $(x-3)^2 + (y+2)^2 + (z-1)^2 = 100$ và mặt phẳng $(\alpha)$ có phương trình $2x - 2y - z + 9 = 0$. Tính bán kính của đường tròn $(C)$ là giao tuyến của mặt phẳng $(\alpha)$ và mặt cầu $(S)$.}{8}{1}{Mặt cầu $(S)$ có tâm $I(3;-2;1)$, bán kính $R = 10$.
Khoảng cách $d(I, (\alpha)) = \frac{|2.3 - 2(-2) - 1 + 9|}{\sqrt{2^2 + (-2)^2 + (-1)^2}} = \frac{18}{3} = 6$.
Bán kính đường tròn giao tuyến $r = \sqrt{R^2 - d^2} = \sqrt{100 - 36} = 8$.}

\shortanswer{Trong không gian $Oxyz$, cho mặt phẳng $(P): 2x - y - 2z - 1 = 0$ và điểm $M(1;-2;0)$. Mặt cầu tâm $M$, bán kính bằng $\sqrt{3}$ cắt phẳng $(P)$ theo giao tuyến là đường tròn có bán kính bằng bao nhiêu?}{1{,}4}{1}{Khoảng cách $d(M, (P)) = \frac{|2.1 - (-2) - 2.0 - 1|}{\sqrt{2^2 + (-1)^2 + (-2)^2}} = \frac{3}{3} = 1$.
Bán kính đường tròn giao tuyến là $r = \sqrt{R^2 - d^2} = \sqrt{(\sqrt{3})^2 - 1^2} = \sqrt{2} \approx 1{,}4$.}

\shortanswer{Trong không gian $Oxyz$, cho mặt phẳng $(P): x + 2y - 2z + 3 = 0$ và mặt cầu $(S)$ có tâm $I(0;-2;1)$. Biết mặt phẳng $(P)$ cắt mặt cầu $(S)$ theo giao tuyến là một đường tròn có diện tích $2\pi$. Một đường thẳng bất kì cắt mặt cầu $(S)$ tại hai điểm $A, B$ sao cho $AB$ lớn nhất. Tính độ dài đoạn thẳng $AB$.}{3{,}5}{1}{Diện tích đường tròn $S = \pi r^2 = 2\pi \Rightarrow r^2 = 2$.
Khoảng cách $d(I, (P)) = \frac{|0 + 2(-2) - 2.1 + 3|}{\sqrt{1^2 + 2^2 + (-2)^2}} = \frac{3}{3} = 1$.
Bán kính mặt cầu $R = \sqrt{r^2 + d^2} = \sqrt{2 + 1} = \sqrt{3}$.
Độ dài $AB$ lớn nhất khi $AB$ là đường kính của mặt cầu $\Rightarrow AB = 2R = 2\sqrt{3} \approx 3{,}5$.}

\shortanswer{Trong không gian $Oxyz$, cho bốn mặt cầu có bán kính lần lượt là $2, 3, 3, 2$ (đơn vị độ dài) tiếp xúc ngoài với nhau. Mặt cầu nhỏ nhất tiếp xúc ngoài với cả bốn mặt cầu nói trên có bán kính bằng bao nhiêu?}{0{,}5}{1}{Áp dụng định lý Soddy cho 5 mặt cầu tiếp xúc ngoài nhau với độ cong $k_i = \frac{1}{r_i}$:
$k_1 = \frac{1}{2}, k_2 = \frac{1}{3}, k_3 = \frac{1}{3}, k_4 = \frac{1}{2}$.
Ta có $S = \sum_{i=1}^4 k_i = \frac{5}{3}$ và $Q = \sum_{i=1}^4 k_i^2 = \frac{13}{18}$.
Phương trình xác định độ cong $k_5$: $2k_5^2 - 2S k_5 + (3Q - S^2) = 0 \Leftrightarrow 36k_5^2 - 60k_5 - 11 = 0 \Rightarrow k_5 = \frac{11}{6}$.
Bán kính mặt cầu nhỏ nhất: $r_5 = \frac{1}{k_5} = \frac{6}{11} \approx 0{,}5$.}

\shortanswer{Trong không gian với hệ toạ độ $Oxyz$, cho điểm $A(1;-2;3)$ và đường thẳng $d$ có phương trình $3x - 2y + z - 7 = 0$. Bán kính mặt cầu tâm $A$, tiếp xúc với $d$ bằng bao nhiêu?}{0{,}8}{1}{Bán kính mặt cầu $R = d(A, d) = \frac{|3.1 - 2(-2) + 3 - 7|}{\sqrt{3^2 + (-2)^2 + 1^2}} = \frac{3}{\sqrt{14}} = \frac{3\sqrt{14}}{14} \approx 0{,}8$.}

\shortanswer{Trong không gian $Oxyz$, cho đường thẳng $d: \frac{x+5}{2} = \frac{y-7}{-2} = \frac{z}{1}$ và điểm $M(4;1;6)$. Đường thẳng $d$ cắt mặt cầu $(S)$ có tâm $M$ tại hai điểm $A, B$ sao cho $AB = 6$. Đường kính của mặt cầu $(S)$ bằng bao nhiêu?}{8{,}5}{1}{Đường thẳng $d$ đi qua $N(-5;7;0)$ và có VTCP $\vec{u}=(2;-2;1)$.
Ta có $\vec{MN} = (-9;6;-6) \Rightarrow [\vec{MN},\vec{u}] = (-6;-3;6)$.
Khoảng cách từ $M$ đến $d$: $h = d(M,d) = \frac{\sqrt{(-6)^2 + (-3)^2 + 6^2}}{\sqrt{2^2 + (-2)^2 + 1^2}} = \frac{9}{3} = 3$.
Bán kính mặt cầu $R = \sqrt{h^2 + \left(\frac{AB}{2}\right)^2} = \sqrt{3^2 + 3^2} = 3\sqrt{2}$.
Đường kính mặt cầu: $D = 2R = 6\sqrt{2} \approx 8{,}5$.}

\end{quiz}

\end{document}
